% Flavor catalog per-process template.
% process_id: CR010
% family: collider_rs
% owner_agent_id: PKA-CR010
% writer_agent_id: null
% checker_agent_id: null
% opus_signoff_id: null
% Canonical metadata sidecar: CR010.yaml
% Status is controlled by the YAML sidecar status_history, not this comment.

\section{\texorpdfstring{Pair Production of a Vector-Like \((T,B)\) Doublet}{Pair Production of a Vector-Like (T,B) Doublet}}

\subsection*{Process}
This entry covers direct LHC searches for QCD pair production of a
vector-like quark doublet,
\[
  pp \to (T,B)(\bar T,\bar B) \to
  (Wb,Zt,Ht)(Wt,Zb,Hb),
\]
with \(T\) carrying charge \(+2/3\) and \(B\) carrying charge \(-1/3\).
The experimental analyses are usually reported as pair production of
\(T\bar T\) and \(B\bar B\), with limits then interpreted in singlet or
doublet branching-fraction scenarios.  The CR010 scope is the weak-isospin
\((T,B)\) doublet and its mixed third-generation final states.

\subsection*{Current best limit(s)}
The canonical doublet-specific bound is the ATLAS Run-2 combination recorded
in the PDG 2025 \(t'\) and \(b'\) listings:
\[
  m_T > 1.37~\mathrm{TeV}, \qquad
  m_B > 1.37~\mathrm{TeV}
  \quad (95\%~\mathrm{CL})
\]
for a weak-isospin \((T,B)\) doublet with \(|V_{Tb}|\ll |V_{tB}|\).  The
same local PDG extract is stored as
\texttt{flavor\_catalog/references/CR010/pdg2025\_tprime\_bprime\_vlq\_limits.txt}
and the strict metadata live in \texttt{CR010.yaml}.

Newer full-Run-2 searches set stronger bounds in simplified single-branching
scenarios, but those are not simultaneous \((T,B)\)-doublet exclusions.  The
PDG 2025 \(t'\) listing quotes \(m_T>1.70~\mathrm{TeV}\) for
\(\mathcal B(T\to Wb)=1\) from ATLAS, while the \(b'\) listing quotes CMS
limits \(m_B>1.57~\mathrm{TeV}\), \(1.56~\mathrm{TeV}\), and
\(1.54~\mathrm{TeV}\) for \(\mathcal B(B\to Hb)=1\),
\(\mathcal B(B\to Wt)=1\), and \(\mathcal B(B\to Zb)=1\), respectively.
These simplified endpoints are useful reach markers; the \(1.37~\mathrm{TeV}\)
ATLAS combination remains the cleanest current number for the explicit
\((T,B)\) doublet assignment.

\subsection*{Relevance to RS / anarchic-flavor pipeline}
Custodial Randall--Sundrum models contain heavy vector-like fermion partners
of the third generation, and the direct searches constrain those partner
masses whenever the spectrum, charges, and prompt decays resemble the LHC
simplified models.  The most robust piece is the QCD pair-production rate;
the interpretation still depends on whether the first custodial partners are
identified with the searched \(T\) and \(B\), whether they are approximately
mass degenerate, and which of \(Wb,Zt,Ht,Wt,Zb,Hb\) saturate the branching
fractions.

For the anarchic-flavor scan in this repo, CR010 is a cross-check rather than
the leading mass setter.  The quark-scan methodology note reports
\(\Mkk^{\min}(p50,\gs=3)=47.26~\mathrm{TeV}\) for the low-energy
\(\Delta F=2\) envelope, far above the \(1\)--\(2~\mathrm{TeV}\) direct VLQ
partner reach.  Collider VLQ searches therefore matter most for non-anarchic
or custodial spectra in which the low-energy flavor pressure is softened, or
for diagnosing whether a nominally allowed point predicts third-generation
partners that should already have been seen.

\subsection*{Post-2010 developments}
Early reinterpretation work after the theoretical VLQ search taxonomy showed
how limits depend on branching fractions rather than only on a single decay
mode.  Garberson--Golling used CMS multilepton data to quote a
\((T,B)\)-doublet exclusion around \(640~\mathrm{GeV}\), establishing the
branching-fraction-plane logic used by later searches.  CMS then moved to
13 TeV boosted-object searches; the 2017 single-lepton analysis excluded
doublet-scenario \(T\) masses below \(830~\mathrm{GeV}\).

The major Run-2 consolidation was the ATLAS 2018 combination
\texttt{arXiv:1808.02343}, which combined \(T\bar T\) and \(B\bar B\) channels
and set the current doublet-specific \(1.37~\mathrm{TeV}\) mass limit.  Full
Run-2 analyses then improved simplified-mode reach: CMS
\texttt{arXiv:2209.07327} used single-lepton, same-sign dilepton, and
multilepton final states with \(138~\mathrm{fb}^{-1}\), and ATLAS
\texttt{arXiv:2401.17165} pushed the pure \(T\to Wb\) endpoint to
\(1.70~\mathrm{TeV}\).  The CMS 2025 review \texttt{arXiv:2405.17605}
summarizes the Run-2 VLQ search program and records the move toward combined
and projection studies.

\subsection*{Constraint validity / model dependence}
The published pair-production limits assume prompt decays, narrow or
analysis-compatible resonance widths, and decay topologies dominated by
third-generation Standard Model bosons plus quarks.  The doublet limit also
assumes the weak-isospin \((T,B)\) pattern and \(|V_{Tb}|\ll |V_{tB}|\); it is
not a model-independent lower bound on every custodial-RS fermion.  Cascades
through other KK states, compressed spectra, nonstandard decays, or large
single-production contributions require a dedicated reinterpretation.

For RS use, the collider observable is therefore best treated as a mass-limit
surface conditional on spectrum and branching-fraction inputs, not as a
single universal \(M_{KK}\) cut.  It constrains partner masses directly; only
a model-specific spectrum calculation can translate that into a bound on the
gauge KK scale or the flavor scan's \(\Mkk\).

\subsection*{Code coverage in this repo}
\textbf{NO}.  A refined grep for VLQ, vector-like, top-partner, bottom-partner,
ATLAS/CMS collider-limit, and reinterpretation-tool keywords found no CR010
implementation in \texttt{quarkConstraints/}, \texttt{qcd/},
\texttt{flavorConstraints/}, \texttt{neutrinos/}, \texttt{yukawa/},
\texttt{warpConfig/}, \texttt{solvers/}, \texttt{scanParams/}, or
\texttt{tests/}.  The only ATLAS/CMS-adjacent grep hit was
\texttt{tests/test\_alpha\_s.py:89}, a CMS/RunDec \(\alpha_s\) example, not a
collider constraint.

Adjacent evidence confirms the gap.  \texttt{quarkConstraints/scan.py:359--441}
computes \(\Mkk\), evaluates \(\Delta F=2\) constraints, and records
\(\epsilon_K\), \(B_d\), \(B_s\), and \(D\)-mixing pass/fail ratios, with no
ATLAS/CMS mass-exclusion gate.  \texttt{quarkConstraints/validation.py:108--145}
uses the same \(\Delta F=2\) summary for benchmark validation.

\subsection*{Implementation difficulty}
\textbf{HIGH}.  A useful live CR010 constraint would need a spectrum-to-final-
state map, branching fractions, efficiencies or simplified-likelihood inputs,
and a recast layer such as CheckMATE, MadAnalysis5, or SModelS.  A hard-coded
mass threshold would be misleading because the published exclusion changes
substantially across branching fractions, mass splittings, widths, and
nonstandard RS decay chains.

\subsection*{Key references}
\texttt{pdg2025\_tprime\_bprime\_vlq\_limits};
\texttt{aguilar\_saavedra\_2009\_arxiv0907\_3155};
\texttt{garberson\_golling\_2013\_arxiv1301\_4454};
\texttt{cms\_2017\_arxiv1706\_03408};
\texttt{atlas\_2018\_arxiv1808\_02343};
\texttt{cms\_2023\_arxiv2209\_07327};
\texttt{atlas\_2024\_arxiv2401\_17165};
\texttt{cms\_review\_2025\_arxiv2405\_17605}.
